% Section 1.4, from lectures/L01/appendix/c_exercises.md.

\section{Exercises}
\label{c1:sec:exercises}\label{c1:app:c}

\begin{aside}
Eight, ending with the Cross-check. Do them in order; the last one needs the script from
Exercise~\ref{c1:ex:6}.

Several of these have a plausible wrong answer that is worth walking into before you find it. Where
an exercise can be checked mechanically, a \textbf{Check yourself} line says how.

Everything here runs against the trees \code{make kernel} and \code{make rootfs} already downloaded,
under \file{build/}. Nothing needs the network.

\textbf{Solutions.} The exercises that are not Code are answered in \secref{sol:sec:c1}. The Code
exercises have no solution: each is checked by running it, as its Check yourself line says.
\end{aside}

% ------------------------------------------------------------------------------------------
\exercise[fillaccent2]{The four pieces}{Recall}
\label{c1:ex:1}

Without looking at \S\ref{c1:app:a}, name the four pieces of an embedded Linux system and say in one
sentence what each does.

Then, for each of the following, say which piece it belongs to:

\xpart{a} \file{/lib/modules/6.12.30/kernel/lib/kunit/kunit.ko}

\xpart{b} \file{u-boot.img}

\xpart{c} \code{aarch64-linux-gnu-gcc}

\xpart{d} \file{/init}

\xpart{e} The \code{.dtb} handed to the kernel at boot

\xpart{f} \file{/bin/busybox}

% ------------------------------------------------------------------------------------------
\exercise[fillaccent2]{Why there are two bootloaders}{Recall}
\label{c1:ex:2}

A colleague looking at a board's flash layout asks why there is both an \code{SPL} and a
\file{u-boot.img}, and suggests deleting the first to save space.

\xpart{a} Explain in two sentences why the SPL exists.

\xpart{b} What specifically would fail if you deleted it, and at what point in the boot?

\xpart{c} Name one thing about a board that makes bricking it recoverable, and one that does not.

% ------------------------------------------------------------------------------------------
\exercise[fillaccent3]{What the toolchain triplet promises}{Hand calculation}
\label{c1:ex:3}

\xpart{a} Break \code{aarch64-linux-gnu} into its parts and say what each names.

\xpart{b} You are handed a binary and told only that it was built with \code{arm-none-eabi-gcc}.
What can you say about the system it is intended to run on, and in particular, can it run on Linux?
Justify it from the triplet alone.

\xpart{c} A program builds against \code{aarch64-linux-gnu-gcc} on a machine whose sysroot has glibc
2.36. It is deployed to a target whose root filesystem has glibc 2.31. Predict what happens, and say
whether the failure is at build time or at run time.

\xpart{d} A colleague fixes a ``missing header'' error in a cross build by running \code{sudo apt
install libfoo-dev}. Explain why this does not work, and what would.

% ------------------------------------------------------------------------------------------
\exercise[fillaccent4]{What leaves the building}{Design}
\label{c1:ex:4}

You ship a product containing exactly these things:

\begin{booktable}{@{}>{\hsize=1.06\hsize}Ll>{\hsize=0.94\hsize}L@{}}
  \thead{Component} & \thead{Licence} & \thead{Modified by you} \\
  \midrule
  U-Boot & GPL-2.0-or-later & Yes, a board port \\
  Linux kernel & GPL-2.0-only & Yes, one driver in-tree \\
  BusyBox & GPL-2.0-only & No \\
  \code{libcurl} & MIT-like (curl) & No \\
  A logging daemon you wrote & Proprietary & n/a \\
  A kernel module you wrote & You choose & n/a \\
\end{booktable}

\noindent A customer sends a written request for the source code.

\xpart{a} Go through the table and say, for each row, what you must publish.

\xpart{b} BusyBox is unmodified. Does that change the answer? Say why.

\xpart{c} Your logging daemon links dynamically against \code{libcurl} and talks to your kernel
module through \code{ioctl}. Does either of those relationships oblige you to publish the daemon?

\xpart{d} What must you keep, and for how long, in order to be able to answer this request at all?

% ------------------------------------------------------------------------------------------
\exercise[fillaccent4]{Where the boundary falls}{Design}
\label{c1:ex:5}

For each pair below, say whether the two things are one work or two, and name the argument that
decides it. Where the answer is genuinely contested, say so rather than picking.

\xpart{a} A proprietary application and the kernel, communicating by system calls.

\xpart{b} A proprietary program statically linked against a GPL-2.0 library.

\xpart{c} A proprietary program dynamically linked against an LGPL-2.1 library.

\xpart{d} A proprietary program dynamically linked against a GPL-2.0 library.

\xpart{e} A proprietary kernel module and the kernel.

\xpart{f} A GPL program and a proprietary program on the same filesystem image, which never
communicate.

% ------------------------------------------------------------------------------------------
\exercise{An SPDX inventory}{Code}
\label{c1:ex:6}

Write a script, \file{lectures/L01/lab/spdx-audit.sh}, that takes a directory and reports what
licences the files under it are under.

It should:
\begin{enumerate}[label=\textbf{\alph*)}]
  \item Walk every \code{*.c} and \code{*.h} file under the directory it is given.
  \item For each, extract the SPDX identifier if there is one. The tag looks like
    \code{SPDX-License-Identifier: GPL-2.0-only} and appears in a comment in roughly the first few
    lines.
  \item Print a count per identifier, most common first.
  \item Print, separately and by name, every file that has no identifier at all. These are the
    interesting ones and they must not be silently folded into a total.
  \item Print the total number of files examined, so that the counts can be checked against it.
\end{enumerate}
Two requirements that are easy to miss and are the point of the exercise:
\begin{itemize}
  \item \textbf{Normalise the deprecated spellings.} SPDX renamed several identifiers;
    \code{GPL-2.0} is the old spelling of \code{GPL-2.0-only}, and \code{GPL-2.0+} is the old
    spelling of \code{GPL-2.0-or-later}. Both spellings are in the kernel tree, in quantity. A
    report that lists them separately is reporting two licences where there is one. Your script must
    report the current spelling and say how many of each it merged.
  \item \textbf{Do not stop at the first match in the file.} Some files carry more than one SPDX
    tag. Take the first one, but know that you have made that choice.
\end{itemize}

\checkyourself run it on \file{build/linux-6.12.30/drivers/rtc}, which has 184 \code{.c} and
\code{.h} files at its top level. Your total must be 184.

% ------------------------------------------------------------------------------------------
\exercise{Measure the machine you built}{Code}
\label{c1:ex:7}

Answer each of these with a command, run on the target after \code{make boot}, or on the host
against \file{build/}. Write down the command as well as the number.

\xpart{a} What kernel version is running, what compiler built it, and when?

\xpart{b} What did the bootloader pass to the kernel on its command line?

\xpart{c} How large is the kernel image, and how large is the compressed initramfs?

\xpart{d} How many modules did the kernel build, and how many are installed on the target?

\xpart{e} How much memory does the target have, and how much of it is free once it has booted?

\xpart{f} What is the physical address of the \code{qa-dev} registers, according to
\file{/proc/iomem}?

For \textbf{f)}, note that nothing has claimed the region yet, because no driver for the device
exists until Chapter~\ref{c6:ch}. Say what you can see and what you cannot.

% ------------------------------------------------------------------------------------------
\exercise[fillaccent]{An audit by hand against an audit by script}{Cross-check}
\label{c1:ex:8}

This is the exercise the other seven exist for. You will compute a set of numbers by hand, run your
own code on the same problem, and reconcile the two. They will not agree, and the disagreement is
the lesson.

\xpart{a} By hand, with no scripting, classify the licences of these ten files in
\file{build/linux-6.12.30/drivers/rtc}:

\begin{console}
rtc-ds1307.c   rtc-pl031.c    rtc-starfire.c   rtc-rs5c313.c   rtc-cmos.c
rtc-m41t80.c   rtc-pcf8563.c  rtc-abx80x.c     rtc-max77686.c  rtc-s35390a.c
\end{console}

\noindent For each, write down the identifier and where you found it. Do not run your script yet.

\xpart{b} Now predict, from that sample of ten, what the breakdown over the whole of
\file{drivers/rtc} will be. Write down a number of distinct licences and a rough proportion for
each.

\xpart{c} Run your script from Exercise~\ref{c1:ex:6} over the whole directory. Record the output.

\xpart{d} Reconcile. Specifically, answer these:
\begin{itemize}
  \item How many \emph{distinct SPDX strings} did your script find, and how many \emph{distinct
    licences} is that actually? What is the relationship between the two numbers?
  \item Your script reports some number of untagged files. Find each one and determine its licence
    by another route. Say what that route was, and whether it was the same route for both.
  \item Did your ten-file sample predict the whole? If your proportions were off, say which
    direction and why a sample drawn from files you recognised the names of would be biased.
\end{itemize}

\xpart{e} A colleague runs your script, sees no \code{Proprietary} in the output, and concludes that
the kernel tree contains nothing proprietary. Say what is wrong with that inference. There are at
least two separate problems with it.
